\section{Discussion and Extensions}

The results matter most in empirical settings where the multiplicative conditional mean is substantively natural and the treatment timing is staggered. Gravity applications are a leading example: PPML specifications with high-dimensional fixed effects are standard because trade flows are modeled through multiplicative resistance terms, exposure variables, and policy shifters, and because zeros are common in the outcome \citep{AndersonVanWincoop2003,HeadMayer2014,YotovPiermartiniMonteiroLarch2016,SantosSilvaTenreyro2006,SantosSilvaTenreyro2011,CorreiaGuimaraesZylkin2020}. In such designs, a pooled treatment indicator may look like a compact summary of a policy change. The results above show that, under heterogeneous proportional effects, the population coefficient of that regression is instead a fixed-effect Poisson projection whose sign can be governed by residualized comparisons across cohort-time cells.

The conclusion pinpoints how PPML summarizes a multiplicative mean model. The homogeneous benchmark in \cref{obj:prop:homogeneous-effect-reduction} shows that a common proportional effect is recovered exactly. With heterogeneous cohort-time proportional effects, \cref{obj:thm:sharp-ppml-forbidden-sign} shows that increasing a positive treated-cell effect moves the pooled coefficient in the direction determined by the weighted fixed-effect residual for that cell. The interpretation problem therefore arises from compressing a collection of heterogeneous economic effects into a single pooled projection coordinate.

\subsection{Limitations and future work}

The primitive index \(\Phi\) in \cref{obj:thm:primitive-global-frontier} is a population diagnostic conditional on known or specified primitives. It determines whether the limiting fixed-effect PPML treatment coordinate is negative, zero, or positive. Estimation of those primitives, implementation of \(\Phi\) as a specification test, and construction of positive-weight estimators from granular proportional effects are the main next steps.

That distinction is important for empirical interpretation. A negative pooled PPML coefficient is sometimes read as evidence that the treatment reduced the conditional mean. Under the conditions of \cref{obj:thm:primitive-global-frontier}, the same primitive configuration can instead have \(\beta^\star(\delta)<0\) and a positive counterfactual-share proportional treatment-on-the-treated target. The sign reversal therefore identifies the pooled fixed-effect PPML coefficient as a projection summary distinct from positive-weight proportional targets.

The positive counterfactual-share target remains closer to the causal object under the maintained proportional-effect restrictions because it aggregates treated-cell effects with weights tied to untreated counterfactual exposure \citep{MoreauKastler2025PTT}. When each granular proportional effect is positive, such an aggregation preserves the sign by construction. \Cref{obj:thm:primitive-global-frontier} formalizes this contrast in the same population environment: the pooled coefficient can be negative while \(PTT\) is positive. Empirical conclusions should therefore distinguish the regression projection from the proportional treatment-on-the-treated estimand.

The four-cohort witness in \cref{obj:prop:four-cohort-sign-reversal} establishes the sign reversal with equal cohort shares, flat untreated time effects, and positive treated-cell effects. Staggered timing and heterogeneous proportional effects place a large positive effect in a cell with a negative residualized treatment value. This makes the phenomenon relevant for conventional staggered-adoption designs, including policy settings where treatment intensity or exposure plausibly varies across treated cohorts and periods.

The population characterization also clarifies the choice that precedes estimation. A researcher may intend a common log multiplier, a collection of granular proportional effects, or a positive-weight proportional aggregate; the pooled projection equals the first under homogeneity and can differ sharply from the third under heterogeneity. The primitive index \(\Phi\) establishes this distinction theoretically. Future work can turn that diagnostic into an operational procedure by estimating granular effects, constructing a sampling distribution for \(\Phi\), assessing separation, and determining how often the sign-reversal region occurs in empirically calibrated designs \citep{SantosSilvaTenreyro2006,SantosSilvaTenreyro2011,CorreiaGuimaraesZylkin2020,Wooldridge2023NonlinearDiD,NagengastYotov2025}.
